\documentclass{amsart}
\usepackage{amsmath,amssymb,amsthm}
% !LOOM begin loom-macros
% loom.sty is MIT-licensed so that it can travel with a paper; see the LICENSE beside it in the loom distribution.

% Dependencies the text does not express as \ref or a matchable \cite.
% Prints nothing.
\providecommand{\uses}[1]{}

% Marks a statement or proof as mathematically incomplete.
% Prints nothing; an author may \renewcommand it to print in draft builds.
\providecommand{\incomplete}[1]{}

% Inputs a file with every sectioning command inside it shifted one level down.
% Composes when nested files nest files.
\providecommand{\nest}[1]{\begingroup
  \let\chapter\section
  \let\section\subsection
  \let\subsection\subsubsection
  \let\subsubsection\paragraph
  \let\paragraph\subparagraph
  \input{#1}\endgroup}
% !LOOM end loom-macros

\theoremstyle{plain}
\newtheorem{theorem}{Theorem}[section]
\newtheorem{lemma}[theorem]{Lemma}
\newtheorem{proposition}[theorem]{Proposition}
\theoremstyle{definition}
\newtheorem{definition}[theorem]{Definition}
\theoremstyle{remark}
\newtheorem{remark}[theorem]{Remark}

\newcommand{\Fix}{\operatorname{Fix}}

\title{Widgets and their fixed loci}
\author{The loom demo}

\begin{document}
\maketitle

\section{Setup}\label{dm-0010}

A widget is the simplest object with an involution. We fix conventions here and prove the two results that every later section uses.

% !LOOM author: The loom demo
% !LOOM created: 2026-09-16
% !LOOM tags: setup

\begin{definition}[Widget]\label{dm-0001}\label{def:widget}
A \emph{widget} is a pair $(X,\sigma)$ of a set $X$ and a map $\sigma\colon X\to X$ with $\sigma^2=\mathrm{id}$. Its \emph{fixed locus} is
\begin{equation}\label{eq:fix}
\Fix(\sigma) = \{x\in X : \sigma x = x\}.
\end{equation}
\end{definition}


% !LOOM author: The loom demo
% !LOOM created: 2026-09-16
% !LOOM tags: setup, orbits

\begin{lemma}[Orbits]\label{dm-0002}\label{lem:orbits}
Every orbit of a widget has one or two points, and $\Fix(\sigma)$ is the union of the one-point orbits.
\end{lemma}
\begin{proof}
Since $\sigma^2=\mathrm{id}$, the orbit of $x$ is $\{x,\sigma x\}$, which has one point exactly when $\sigma x = x$; compare \eqref{eq:fix}.
\end{proof}


\section{Results}\label{dm-0011}

The main theorem says that the fixed locus is as nice as one could hope.

% !LOOM author: The loom demo
% !LOOM created: 2026-09-16
% !LOOM tags: main

\begin{theorem}[Main]\label{dm-0003}
Let $(X,\sigma)$ be a widget with $X$ a finite set. Then $\Fix(\sigma)$ is nonempty if and only if $|X|$ is odd, and it is closed in every topology on $X$ for which $\sigma$ is continuous and $X$ is Hausdorff.
\end{theorem}
\begin{proof}
\uses{dm-0002}
By Lemma~\ref{lem:orbits} the set $X$ is a disjoint union of orbits of size one or two, so $|X| \equiv |\Fix(\sigma)| \pmod 2$, which gives the parity statement. For closedness, $\Fix(\sigma)$ is the preimage of the diagonal under $(\mathrm{id},\sigma)\colon X\to X\times X$, and the diagonal is closed when $X$ is Hausdorff; this is \cite[Proposition 3.2]{Calloway14}.
\end{proof}


\begin{remark}\label{dm-0004}
Closedness is where the involution being an automorphism matters; a merely continuous $\sigma$ can have a fixed locus that is not closed in the constructible topology.
\end{remark}

\begin{proposition}[Functoriality]\label{dm-0005}
A map of widgets $f\colon (X,\sigma)\to(Y,\tau)$ with $f\sigma = \tau f$ restricts to a map $\Fix(\sigma)\to\Fix(\tau)$.
\end{proposition}
\begin{proof}
\uses{dm-0001}
If $\sigma x = x$ then $\tau f(x) = f(\sigma x) = f(x)$. \incomplete{Say why the restriction is continuous when $X$ carries the constructible topology.}
\end{proof}

\bibliographystyle{amsplain}
\bibliography{refs}

\end{document}
